\documentclass[11pt,a4paper]{article}
\usepackage[margin=2.5cm]{geometry}
\usepackage{fancyhdr}
\usepackage{parskip}
\usepackage[utf8]{inputenc}
\usepackage[T1]{fontenc}

\pagestyle{fancy}
\fancyhf{}
\fancyhead[L]{\small Exam Answers}
\fancyhead[R]{\small The Culture of Business}
\fancyfoot[C]{\thepage}

\title{Second Exam Answers\\
\large The Culture of Business, Finance and the Economy}
\author{Juan Ignacio Elosegui\\
\small Universidad Torcuato Di Tella\\
\small 2026}
\date{}

\begin{document}
\maketitle

\textbf{Instructions:} Answer \textbf{4 questions} in total: at least 3 from questions 1--6 and no more than 1 from questions 7--9. This document prepares 4 answers: \textbf{Q1} (option B), \textbf{Q2}, \textbf{Q3} (mandatory group) and \textbf{Q7} (elective group).

\section*{Question 1 --- The Self as Business (Class XVII)}

\textbf{Question:} \textit{Discussing its emergence and characteristics, characterise the ``self as business'' and show its workings in TCBFE as thoroughly as possible and with as many class examples as possible.}

To understand how the self became a business, we need to know how the very concept of selfhood evolved in relation to work.
\begin{itemize}
	\item \textbf{Antiquity}. ``Being yourself'' was a matter of property --- owning your house, your land, your possessions.
	\item \textbf{Industrial revolution}. The self became a \textbf{property} that could be rented out to a factory or a firm for a specific purpose. Two features defined this era:
	      \subitem \textbf{(1) Exclusivity}: the worker committed to one employer, and the employer reciprocated with loyalty for an implicit contract. Labour was scarce, so companies encouraged long-term commitment through stable pay, insurance, promotions, and what amounted to \textit{treats as rewards for loyalty}.
	      \subitem \textbf{(2) Boundaries}: work happened from 9 to 5, and you lived your life outside of that boundary. What you did after hours was your own affair.
	\item \textbf{Contemporary.} The self is no longer a property to be rented --- it has become \textbf{a business}. ``I am a bundle of skills, assets and experiences that must be managed and enhanced.'' ``We are the CEOs of our own companies: Me Inc.'' To be employable now means to \textit{represent yourself as a business}.
	      \subitem \textbf{(1) Loyalty was replaced by passion.} Passion is an individual value, loyalty isn't. If you do not have passion, you will not get the job. For passion, companies no longer offer houses or mortgages --- they offer gym memberships and discounts. You are \textit{no longer an employee}; you become a ``\textbf{collaborator}'' because you ``work together'' --- two businesses collaborating with each other. Loyalty matters so little that no exclusivity is expected or offered any more.
	      \subitem \textbf{(2) Boundaries are dead.} Since we are expected to be available on demand and always ``on,'' whatever we do outside working hours affects our professional standing.
\end{itemize}

The self-as-business manifests concretely in four ways:
\begin{itemize}
	\item \textbf{Skills}: historically, skills were concrete and job-related. Now, things all humans do --- like communication or replying to emails --- have been ``assetized.'' Speaking and listening became a ``skill.'' You are expected to know and rank your skills as if they were product features.
	\item \textbf{CVs and LinkedIn}: when we were property, a CV was a historical record of what you had done. Now a resume is no longer that --- it is a \textbf{marketing document}. We no longer tell people what we do; we present metrics and results. Back in the day, references and grades would have done the same job.
	\item \textbf{Personal brand}: if you are your own CEO, you are the brand. Since staying at a company has no value and people switch jobs constantly, people must \textbf{craft a narrative about their work history} for recruiters. The cohesiveness of that narrative is crucial.
	\item \textbf{Networking}: there are far too many jobs and too much demand. Employers increasingly automatize CV reception through \textbf{applicant tracking systems}, which are designed to manage workload, not find the best candidate. Networking becomes crucial to bypass these systems because your contacts know you.
\end{itemize}

Skills and knowledge \textit{are} a form of capital produced by deliberate investment. The body itself becomes an asset: sleep, nutrition, and exercise should be framed not as personal needs but as \textbf{investments} --- ``if you are an important asset, how could depriving and depreciating that asset possibly be good for your organization?'' Even personal trauma and grief become monetizable raw material.

\section*{Question 2 --- From Makerspaces to Venture Capitalism: The Very Beginnings (Class XVIII)}

\textbf{Question:} \textit{Retrace the very beginnings of the journey from makerspaces to venture capitalism as thoroughly as possible and with as many class examples as possible.}

The journey from makerspaces to venture capitalism begins in the early 2000s with a particular kind of space and ethos. Hackerspaces, makerspaces, and hackathons emerged as sites of \textbf{playful experimentation} --- fun and flexible, organized around peer production and community. The promise was a shift \textbf{from being a consumer to being a maker, a producer} --- someone who can make stuff. The defining spirit was ``playing for playing's sake.''

This culture was vibrant but uncommercial. People gathered to learn and share. About ten years later, a recurring topic in the hacker and makerspaces was how to ``\textbf{take things to the next level}'' and to ``\textbf{scale up}.'' This was the pivotal moment: the transition from making for its own sake to making for the market.

The institutional vehicle for this transition was the \textbf{incubator}. Startup incubator programs are structured and intensive initiatives designed to support early-stage startups. They provide entrepreneurs with resources, guidance, and mentorship to help them grow their businesses. There is a competitive application process selecting only the most promising startups, a physical workspace fostering collaboration, access to resources, mentoring and guidance from experienced entrepreneurs and investors, networking opportunities...

What is crucial about incubators, is that they appeared to \textbf{retain the commitment to playful experimentation, fun and flexible work culture, peer production, and community} that many had come to identify with maker and hacker culture. The same enthusiasm, the same ``having fun,'' but the promise was now to connect this culture of tinkering to ``the real world.'' The critical switch into entrepreneurship required learning how to \textbf{derive pleasure not only from hacking machines but from hacking markets}. Incubators, in this sense, are ``start-up schools'' that ``coin \textbf{entrepreneurs first and ideas later}.''

UCHINA ran a hardware acceleration program that promised to move products from concept to mass production in months rather than years. Shenzhen was enriched with a sense of \textbf{speed} --- it promised to teach people how to accelerate. The logic was: by ``accelerating'' the time it takes to bring a new product to market, one increases one's attractiveness to investors.

The program trained teams to attract two audiences: \textbf{capital investors} (through pitches) and a \textbf{potential consumer base} (via crowdfunding platforms). Crucially, the crowdfunding campaigns were \textbf{not actually designed to sell products or raise money}. Their real purpose was to show that the product had a potential market. Investors look for these indicators to minimize their risks. A successful campaign shows the promise of future gains. \textbf{The start-ups were taught to produce not products to be sold, but promises to be invested in.}''

This is where makerspaces end and venture capitalism begins. Incubators and accelerators channel the maker spirit of speed, scale, and magnitude directly towards venture capital --- transforming playful tinkerers into investable entrepreneurs, and prototypes into promissory vehicles for financial return.

\section*{Question 3 --- Hype, Pitching, Leveraging, and Venture Capitalism (Class XIX)}

\textbf{Question:} \textit{Describe the workings of ``hype, pitching and leveraging'' and the main features of venture capitalism as thoroughly as possible and with as many examples from class as possible.}

\textbf{Hype} is not about truth or falsity, but about \textbf{credibility or incredibility}. It is a \textbf{credible, fabricated truth} that entrepreneurs and investors work toward. It becomes a way for imagining the future and working to enact that future.

The example of Incyte and the Human Genome Project illustrates this perfectly: the product of a genome company is the creation of \textit{information} that is a condition of possibility for the creation of a drug. The conditions of its own possibility can only be created through a \textit{vision}. This is why the \textbf{vision of the future creates the conditions of possibility for the existence of the company in the present}. Visions do not even have to be true to sell.

\textbf{Pitching} is the performative art through which hype is delivered to investors.

In the STHLM Tech Meetup in Stockholm, two entrepreneurs pitched their startup to a VC panel, opening with a story about forests and climate change. Their actual product was a social logbook mobile app for naturalists that crowdsourced data for research. The correct pitch, according to VCs, would be: ``I've built a great app that \textbf{people love to use}. It's a digital record book for hikers, birders, whatever --- it has \textbf{social stuff, gamification, it's addicting, people want to pay to use it}. How do I know? Because \textbf{we already have ten thousand users}. And we can use the data to save the planet by giving it to scientists.'' An investor shouted: ``I would give you money for that!''

\textbf{Leveraging} means borrowing money from lenders, banks, or investors --- \textbf{taking on risk in order to do business with resources that one does not own} for the sake of speed and flexibility to take things further. Make it or leave.

As for the main features of \textbf{venture capitalism}: VC investors have a \textbf{responsibility to grow their fund} through investments that will \textbf{scale their valuations exponentially} over the term of investment. Even VCs who care about climate change have little to offer on that matter alone --- what they need is the core of the business, growth potential, and a path to getting their money back. Hype, pitching, and leveraging together constitute an \textit{ethos that seizes our rationality from beyond} --- a logic of credible futures, performed promises, and leveraged relationships that sustains the startup world.

\section*{Question 7 --- Consultants' Abstractions (Class XIII)}

\textbf{Question:} \textit{Discuss: ``The abstractions consultants produce and deal with perpetuate the need for their involvement and deflect responsibility.''}

Consultants benefit from and actively cultivate uncertainty --- a \emph{profitable uncertainty}. It is a \textit{kind of speech} and a \textit{kind of orientation} toward the future. It is a resource to be exploited: it justifies the consultant's presence, generates new engagements, and keeps clients dependent.

There are three mechanisms through which consultants make uncertainty profitable, and each of them produces abstractions that perpetuate involvement and deflect responsibility.
\begin{itemize}
	\item \textbf{Potential worlds}: consultants construct possible future scenarios --- desirable or threatening --- that make the client feel the need for expert guidance. By painting vivid alternative futures, they create urgency and justify action. These scenarios are abstractions --- they do not describe reality but possible realities, and they position the consultant as the one who can navigate between them.
	\item \textbf{Free decisions}: consultants present \textit{options, not suggestions}. This shifts the responsibility to the client while allowing the consultant to not be blamed against failure. If things go wrong, it was the client's choice.
	\item \textbf{Consultant engagement}: the ongoing relationship itself becomes the product. Uncertainty never fully resolves, so the consultant is always needed for the next step. The abstraction of ``engagement'' perpetuates the need for continued involvement by design.
\end{itemize}

``\textbf{The Gap}'' is another powerful abstraction. The word ``gap'' is widely used in consulting, they constantly invoke revenue gaps, performance gaps, knowledge gaps. These are illustrated with charts and shapes comparing an entity to a \textit{potential state}. The gap simultaneously diagnoses a problem and implies a solution. Gaps can be ``bridged,'' ``filled,'' or ``closed'' --- and doing so requires the consultant's expertise. The purpose is to present to the client a potential world that is \textit{just} out of reach, perpetuating the need for a consultant that can close it.

The abstractions consultants produce form a self-reinforcing system. They create a language in which problems can never be fully solved (because they are defined in unmeasurable terms), solutions always require further consulting (because uncertainty never resolves), and responsibility for failure always lies elsewhere (because the consultant only offered options, not directions). The abstractions are not a by-product of consulting work; they are its core technology.

\end{document}
